\section{模型评价与推广}

本文四问始终围绕同一微网物理结构逐步扩展：问题一建立基础电能流与储能状态关系，
问题二把全天计划与实际执行分开，问题三增加日内新预报到达后的再决策，问题四把实时电价
纳入随机量。这样处理的直接好处是前问的变量与约束可被后问继承，新增变量各自对应一条题意，
模型之间不存在割裂。

基础模型把外网购电分为直接供负荷与用于储能充电、把富余光伏分为储能充电与弃光，
并在问题二以后用 \(W\) 单独描述已购未用电，因此能清楚追踪外网、光伏与储能之间的流向。
问题二以后只把真正需要提前提交的计划购电作为事前变量，充放电、弃光与紧急购电留待情景实现后决定，
符合真实信息时序；用历史完整日预测误差构造情景，避免了把单一预测曲线当作真值，7日滚动视野又使日末储电量具有未来价值，减轻了单日模型在时域末端过度放电的问题。
问题三用四个时点的新预报重新调整尚未执行的计划，并以真实执行后的储电量连接各阶段，
使“新信息—再决策—实际执行”构成连续链条。问题四保留负价格情景而不截断，
通过计划购电物理上界与调增调减互斥约束消除非物理套利，在不改变原结算规则的前提下保证模型有界。

模型仍存在局限。目标函数只计购电费用，未显式加入储能循环衰减与寿命损耗，
频繁充放电的长期成本可能被低估。正式随机模型的情景数受全年滚动计算量限制，
问题二的结果已表明紧急购电量对情景覆盖较敏感；问题三只更新光伏预报，负荷中心预测仍沿用0:00版本；
问题四的调度效果受价格峰谷时刻预测能力制约。此外，按题面设定未设余电上网与需求响应机制，
适用范围限于“负荷—光伏—储能—外网购电”这一类微网。

后续可在目标函数中加入与充放电通量及循环深度相关的衰减成本，使短期电费与长期寿命同时进入优化；
可增加情景数并配合情景削减控制计算规模；也可把当前的阶段滚动随机规划扩展为严格多阶段随机规划，
使每阶段决策只依赖当时已观察到的信息。基础物理层由负荷、分布式电源、储能与外网组成，
重新给定容量、功率、效率、电价与误差数据后，同一思路可用于园区微网、商业建筑群、
校园能源系统以及含风电等其他随机电源的储能调度。
